\chapter*{Lời Mở Đầu}
\addcontentsline{toc}{chapter}{Lời Mở Đầu}
\markboth{Lời Mở Đầu}{Lời Mở Đầu}

\begin{truyen}{Lời Mở Đầu}{Opening Words}
\chuhoa{C}{ó những câu học sinh nói ra nghe qua tưởng chỉ là cãi cùn. Nhưng nghe kỹ hơn, ta thấy bên dưới thường là một mớ rất thật: áp lực điểm số, nỗi sợ bị coi thường, thói quen trì hoãn, sự mệt mỏi, hoặc đơn giản là không biết nói cho đúng ý mình.}
\textit{(Some student lines sound like mere stubborn excuses. Listen more carefully and you often hear something real underneath: pressure, fear of disrespect, procrastination, exhaustion, or simply poor expression.)}

Vì vậy bản V3 này không cố làm cho mọi cuộc đối thoại trở nên đẹp. Nó giữ nguyên mùi lớp học: có câu vụng, có câu chát, có câu chống chế nghe phát biết ngay là để cứu sĩ diện. Nhưng chính chỗ đó mới là nơi người dạy cần hiểu, và người học cần soi lại.
\textit{(That is why this V3 does not try to beautify every exchange. It keeps the smell of the classroom: awkward lines, sharp replies, and obvious face-saving excuses. That is exactly where teachers need insight and students need reflection.)}

Mỗi đoạn trong cuốn này đều ngắn. Không dài dòng, không giảng cho nhiều. Chỉ là một tình huống, một câu đáp, một cái khựng lại, rồi rút ra điều cốt lõi. Vì ngoài đời, nhiều bài học đến từ đúng một câu bị nói trật chỗ.
\textit{(Each piece is short. No long lectures. Just one situation, one answer, one pause, and one essential takeaway. In real life, many lessons begin with one sentence spoken at the wrong time.)}
\end{truyen}

\begin{baihoc}
Khi lý sự bị bóc trần, điều còn lại không phải là thắng thua, mà là cơ hội để lớn lên.
\textit{(When an excuse is exposed, what remains is not victory or defeat, but a chance to grow.)}
\end{baihoc}
\ngancach